
\documentclass{article}
\usepackage{colortbl}
\usepackage{makecell}
\usepackage{multirow}
\usepackage{supertabular}

\begin{document}

\newcounter{utterance}

\centering \large Interaction Transcript for game `dond', experiment `semi\_en', episode 16 with reflexion{-}Qwen3.5{-}4B.
\vspace{24pt}

{ \footnotesize  \setcounter{utterance}{1}
\setlength{\tabcolsep}{0pt}
\begin{supertabular}{c@{$\;$}|p{.15\linewidth}@{}p{.15\linewidth}p{.15\linewidth}p{.15\linewidth}p{.15\linewidth}p{.15\linewidth}}
    \# & $\;$A & \multicolumn{4}{c}{Game Master} & $\;\:$B\\
    \hline

    \theutterance \stepcounter{utterance}  
    & & \multicolumn{4}{p{0.6\linewidth}}{
        \cellcolor[rgb]{0.9,0.9,0.9}{
            \makecell[{{p{\linewidth}}}]{
                \texttt{\tiny{[P1$\langle$GM]}}
                \texttt{You are playing a negotiation game in which you have to agree on how to divide a set of items among you and another player.} \\
\\ 
\texttt{Rules:} \\
\texttt{(a) You and the other player are given a set of items. Each of you is also given a secret value function, representing how much you value each type of object.} \\
\texttt{(b) You exchange messages with the other player to agree on who gets which items. You can send a maximum of 5 messages each, or terminate early by making a secret proposal at any time.} \\
\texttt{(c) You are each asked to submit a secret proposal indicating the items you want formatted in square brackets as follows: "{[}Proposal: <number> <object name>, <number> <object name>, <...>{]}"} \\
\texttt{(d) If your proposals are complementary, i.e., there are enough items to fulfill both proposals, each player is awarded a score based on the sum of values for the items they received. Otherwise, both of you get zero points.} \\
\texttt{(e) Your goal is to maximize the score you receive.} \\
\\ 
\texttt{Let us start.} \\
\\ 
\texttt{The set of available items is:} \\
\\ 
\texttt{1 paper, 1 egg, 3 onions} \\
\\ 
\texttt{Your secret value function is:} \\
\\ 
\texttt{paper: 4, egg: 0, onion: 2} \\
\\ 
\texttt{IMPORTANT: Your messages, unless it is the secret proposal, are directly transmitted to the other player, so do not include any response to the rules or text announcing your message. To make a secret proposal, use the indicated format. Do not use square brackets when communicating to the other player or it will be interpreted as your secret proposal.} \\
            }
        }
    }
    & \\ \\

    \theutterance \stepcounter{utterance}  
    & \multicolumn{4}{p{0.6\linewidth}}{
        \cellcolor[rgb]{0.9,0.9,0.9}{
            \makecell[{{p{\linewidth}}}]{
                \texttt{\tiny{[P1$\rangle$GM]}}
                \texttt{I value the paper most (4 points), followed by the onions (2 points each), and the egg is worthless to me (0 points). Since I only care about the paper and onions, I should aim to get those. The other player likely values the egg or onions differently.} \\
\\ 
\texttt{To maximize my score, I should propose taking the items I value most. If I take the paper and one onion, I get 6 points. The other player would get the remaining onion and egg.} \\
\\ 
\texttt{Let's try a simple proposal where I take the paper and one onion.} \\
\\ 
\texttt{Paper, 1 onion} \\
            }
        }
    }
    & & \\ \\

    \theutterance \stepcounter{utterance}  
    & & & \multicolumn{2}{p{0.3\linewidth}}{
        \cellcolor[rgb]{0.9,0.9,0.9}{
            \makecell[{{p{\linewidth}}}]{
                \texttt{\tiny{[GM$|$GM]}}
                \texttt{continue} \\
            }
        }
    }
    & & \\ \\

    \theutterance \stepcounter{utterance}  
    & & & \multicolumn{2}{p{0.3\linewidth}}{
        \cellcolor[rgb]{0.9,0.9,0.9}{
            \makecell[{{p{\linewidth}}}]{
                \texttt{\tiny{[GM$|$GM]}}
                \texttt{I value the paper most (4 points), followed by the onions (2 points each), and the egg is worthless to me (0 points). Since I only care about the paper and onions, I should aim to get those. The other player likely values the egg or onions differently.} \\
\\ 
\texttt{To maximize my score, I should propose taking the items I value most. If I take the paper and one onion, I get 6 points. The other player would get the remaining onion and egg.} \\
\\ 
\texttt{Let's try a simple proposal where I take the paper and one onion.} \\
\\ 
\texttt{Paper, 1 onion} \\
            }
        }
    }
    & & \\ \\

    \theutterance \stepcounter{utterance}  
    & & \multicolumn{4}{p{0.6\linewidth}}{
        \cellcolor[rgb]{0.9,0.9,0.9}{
            \makecell[{{p{\linewidth}}}]{
                \texttt{\tiny{[GM$\rangle$P2]}}
                \texttt{You are playing a negotiation game in which you have to agree on how to divide a set of items among you and another player.} \\
\\ 
\texttt{Rules:} \\
\texttt{(a) You and the other player are given a set of items. Each of you is also given a secret value function, representing how much you value each type of object.} \\
\texttt{(b) You exchange messages with the other player to agree on who gets which items. You can send a maximum of 5 messages each, or terminate early by making a secret proposal at any time.} \\
\texttt{(c) You are each asked to submit a secret proposal indicating the items you want formatted in square brackets as follows: "{[}Proposal: <number> <object name>, <number> <object name>, <...>{]}"} \\
\texttt{(d) If your proposals are complementary, i.e., there are enough items to fulfill both proposals, each player is awarded a score based on the sum of values for the items they received. Otherwise, both of you get zero points.} \\
\texttt{(e) Your goal is to maximize the score you receive.} \\
\\ 
\texttt{Let us start.} \\
\\ 
\texttt{The set of available items is:} \\
\\ 
\texttt{1 paper, 1 egg, 3 onions} \\
\\ 
\texttt{Your secret value function is:} \\
\\ 
\texttt{paper: 0, egg: 7, onion: 1} \\
\\ 
\texttt{IMPORTANT: Your messages, unless it is the secret proposal, are directly transmitted to the other player, so do not include any response to the rules or text announcing your message. To make a secret proposal, use the indicated format. Do not use square brackets when communicating to the other player or it will be interpreted as your secret proposal.} \\
\\ 
\\ 
\texttt{I value the paper most (4 points), followed by the onions (2 points each), and the egg is worthless to me (0 points). Since I only care about the paper and onions, I should aim to get those. The other player likely values the egg or onions differently.} \\
\\ 
\texttt{To maximize my score, I should propose taking the items I value most. If I take the paper and one onion, I get 6 points. The other player would get the remaining onion and egg.} \\
\\ 
\texttt{Let's try a simple proposal where I take the paper and one onion.} \\
\\ 
\texttt{Paper, 1 onion} \\
            }
        }
    }
    & \\ \\

    \theutterance \stepcounter{utterance}  
    & & & \multicolumn{4}{p{0.6\linewidth}}{
        \cellcolor[rgb]{0.9,0.9,0.9}{
            \makecell[{{p{\linewidth}}}]{
                \texttt{\tiny{[GM$\langle$P2]}}
                \texttt{Paper, 1 onion} \\
            }
        }
    }
     \\ \\

    \theutterance \stepcounter{utterance}  
    & & & \multicolumn{2}{p{0.3\linewidth}}{
        \cellcolor[rgb]{0.9,0.9,0.9}{
            \makecell[{{p{\linewidth}}}]{
                \texttt{\tiny{[GM$|$GM]}}
                \texttt{continue} \\
            }
        }
    }
    & & \\ \\

    \theutterance \stepcounter{utterance}  
    & & & \multicolumn{2}{p{0.3\linewidth}}{
        \cellcolor[rgb]{0.9,0.9,0.9}{
            \makecell[{{p{\linewidth}}}]{
                \texttt{\tiny{[GM$|$GM]}}
                \texttt{Paper, 1 onion} \\
            }
        }
    }
    & & \\ \\

    \theutterance \stepcounter{utterance}  
    & & \multicolumn{4}{p{0.6\linewidth}}{
        \cellcolor[rgb]{0.9,0.9,0.9}{
            \makecell[{{p{\linewidth}}}]{
                \texttt{\tiny{[P1$\langle$GM]}}
                \texttt{Paper, 1 onion} \\
            }
        }
    }
    & \\ \\

    \theutterance \stepcounter{utterance}  
    & \multicolumn{4}{p{0.6\linewidth}}{
        \cellcolor[rgb]{0.9,0.9,0.9}{
            \makecell[{{p{\linewidth}}}]{
                \texttt{\tiny{[P1$\rangle$GM]}}
                \texttt{I value the paper most (4 points) and the onions (2 points each), while the egg is worthless to me (0 points). To maximize my score, I should secure the paper and at least one onion. If I take the paper and one onion, I get 6 points. The other player would receive the remaining onion and egg. This seems like a fair split that satisfies my high{-}value items.} \\
\\ 
\texttt{{[}Proposal: 1 paper, 1 onion{]}} \\
            }
        }
    }
    & & \\ \\

    \theutterance \stepcounter{utterance}  
    & & & \multicolumn{2}{p{0.3\linewidth}}{
        \cellcolor[rgb]{0.9,0.9,0.9}{
            \makecell[{{p{\linewidth}}}]{
                \texttt{\tiny{[GM$|$GM]}}
                \texttt{proposal} \\
            }
        }
    }
    & & \\ \\

    \theutterance \stepcounter{utterance}  
    & & & \multicolumn{2}{p{0.3\linewidth}}{
        \cellcolor[rgb]{0.9,0.9,0.9}{
            \makecell[{{p{\linewidth}}}]{
                \texttt{\tiny{[GM$|$GM]}}
                [1, 0, 1]
            }
        }
    }
    & & \\ \\

    \theutterance \stepcounter{utterance}  
    & & \multicolumn{4}{p{0.6\linewidth}}{
        \cellcolor[rgb]{0.9,0.9,0.9}{
            \makecell[{{p{\linewidth}}}]{
                \texttt{\tiny{[GM$\rangle$P2]}}
                \texttt{The other player has submitted a secret proposal. It is now your turn to submit a proposal of your own.} \\
            }
        }
    }
    & \\ \\

    \theutterance \stepcounter{utterance}  
    & & & \multicolumn{4}{p{0.6\linewidth}}{
        \cellcolor[rgb]{0.9,0.9,0.9}{
            \makecell[{{p{\linewidth}}}]{
                \texttt{\tiny{[GM$\langle$P2]}}
                \texttt{{[}Proposal: 1 paper, 1 onion{]}} \\
            }
        }
    }
     \\ \\

    \theutterance \stepcounter{utterance}  
    & & & \multicolumn{2}{p{0.3\linewidth}}{
        \cellcolor[rgb]{0.9,0.9,0.9}{
            \makecell[{{p{\linewidth}}}]{
                \texttt{\tiny{[GM$|$GM]}}
                \texttt{proposal} \\
            }
        }
    }
    & & \\ \\

    \theutterance \stepcounter{utterance}  
    & & & \multicolumn{2}{p{0.3\linewidth}}{
        \cellcolor[rgb]{0.9,0.9,0.9}{
            \makecell[{{p{\linewidth}}}]{
                \texttt{\tiny{[GM$|$GM]}}
                [1, 0, 1]
            }
        }
    }
    & & \\ \\

    \theutterance \stepcounter{utterance}  
    & & & \multicolumn{2}{p{0.3\linewidth}}{
        \cellcolor[rgb]{0.9,0.9,0.9}{
            \makecell[{{p{\linewidth}}}]{
                \texttt{\tiny{[GM$|$GM]}}
                [[1, 0, 1], [1, 0, 1]]
            }
        }
    }
    & & \\ \\

    \theutterance \stepcounter{utterance}  
    & \multicolumn{4}{p{0.6\linewidth}}{
        \cellcolor[rgb]{0.95,0.95,0.95}{
            \makecell[{{p{\linewidth}}}]{
                \texttt{\tiny{[P2$\rightarrow$P1]}}
                \texttt{Post{-}episode Reflexion (not sent to GM):} \\
\texttt{New plan: I failed because I made a secret proposal before establishing a mutually beneficial agreement with the other player. In this game, a proposal only yields points if the other player accepts it (implied by "complementary" and the zero{-}score rule for non{-}complementary deals). By proposing immediately, I likely forced a rejection or a suboptimal counter{-}offer without knowing the opponent's valuation. Next time, I will first send a message to the opponent to reveal my strategy or ask for their valuation of the items (specifically the egg, which I value at 0, to see if they value it highly). Once I understand the opponent's priorities, I will craft a proposal that gives them their most valued item in exchange for the items I need, ensuring the proposals are complementary and maximize the sum of our values. For example, if the opponent values the egg highly, I should propose taking the paper and egg (worth 4 to me, but maybe 0 to them) while giving them the onions, or better yet, trade the paper for the onions if the opponent values onions more than the paper. I will only submit a secret proposal after confirming the split is acceptable to the other player through prior negotiation.} \\
            }
        }
    }
    & & \\ \\

    \theutterance \stepcounter{utterance}  
    & & & \multicolumn{4}{p{0.6\linewidth}}{
        \cellcolor[rgb]{0.95,0.95,0.95}{
            \makecell[{{p{\linewidth}}}]{
                \texttt{\tiny{[P1$\rightarrow$P2]}}
                \texttt{Post{-}episode Reflexion (not sent to GM):} \\
\texttt{New plan: I failed because I misread my own value function. The prompt stated I valued paper at 4 points, but I incorrectly assumed I valued it most (4 points) while also saying the egg was worthless (0 points) and onions were 2 points each. However, the critical error was in my strategic logic: I proposed taking the paper and one onion (total value 6) for myself, leaving the other player with the remaining onion (value 1) and the egg (value 7). Since the other player likely values the egg highly (as it's the only high{-}value item left for them if I don't take it), they would reject a deal where they get only 1 point while I get 6, or the system would deem the proposals non{-}complementary if their valuation of the egg is high enough to make their total potential score higher than my offer. I should have analyzed the other player's likely incentives first. Given the items (1 paper, 1 egg, 3 onions) and my values (paper: 4, onion: 2, egg: 0), I should have realized that the egg is worthless to me, so I should not be the one holding it if I can avoid it, but more importantly, I need to ensure the other player gets a deal they accept. If the other player values the egg highly, I should offer them the egg and the paper, and take the onions. My previous proposal gave them the egg and an onion, which might be too little if they value the egg more than the onion. A better strategy is to propose taking the onions (which I value at 2 each) and let the other player take the paper and the egg. This way, I get 2 points (or 6 if I take all 3 onions), and they get the paper (4) and egg (7). I must ensure my proposal explicitly lists the items I want to take, and I should calculate the maximum possible score for myself based on the items I can realistically secure without the other player rejecting the deal. I will propose taking the 3 onions, leaving the paper and egg for the other player.} \\
            }
        }
    }
     \\ \\

\end{supertabular}
}

\end{document}
